\subsection{MPRA数据与固定染色体划分}

我们使用了Gosai等人~\citep{Gosai2024}公开的MPRA序列性能表，包含K562、HepG2和SK-N-SH的活性估计值和标准误。上游染色体划分保持不变：19、21和X号染色体为验证集；7号和13号染色体为测试集；其余符合条件的染色体为训练集。记录要求仅包含A/C/G/T、长度不超过200 nt、含有有限的活性和标准误值、且最大报告log倍变化标准误低于1.0。全数据模型训练前移除精确序列重复。短于200 nt的序列在200位置张量中居中，第五个输入通道标记有效位置。

\subsection{泄漏与亲本重叠审计}

我们检查了200 nt构建体之间的精确同一性、反向互补同一性、主变异组成员关系，以及跨拆分的汉明距离（最多3个）。近重复审计识别出27对跨拆分组对，涉及8个训练侧变异组。亲本拒绝列表包括早期开发和验证亲本，匹配标识符、序列哈希和审计标记的冲突。600亲本主要评估表在不使用活性标签的情况下创建：按\texttt{ID + tab + sequence}的SHA-256对符合条件的200 nt测试序列排序，在拒绝列表移除后选取前600个。消融研究使用固定的180亲本子集，九准则基准使用固定的90亲本验证子集——在发现6个名义亲本与早期探索性基准重叠后，按标识符排除，未参考结果。各阶段的亲本重叠通过标识符和序列哈希检查。

全数据加载器执行了精确去重。独立的汉明距离清单应用于亲本资格审查，确保任何被标记的近重复都不被纳入编辑基准。

\subsection{主预测器}

Gosai等人~\citep{Gosai2024}发布的Malinois BassetBranched作为主预测器。检查点以仅权重模式加载。每条序列在正向和反向互补方向各评估一次，活性预测取平均。对于靶细胞$t$，特异性裕度定义为
\begin{equation}
M_t(x)=f_t(x)-\max_{j\ne t} f_j(x),
\end{equation}
其中$f_j(x)$是细胞类型$j$的预测活性。候选裕度增益为$\Delta M_t=M_t(x_{\mathrm{edit}})-M_t(x_{\mathrm{parent}})$。

\subsection{评审器与跨模型评估集成}

我们训练了两个架构差异化的异方差CNN，在全部质量过滤的训练记录上进行训练，以在不同阶段提供留出评估。评审集成是一个残基-扩张一维CNN，具有11 nt的主干、扩张率1、2、4、8的残基块，以及共享的均值和对数方差头；它在候选搜索期间可用于跨模型重排序。跨模型评估集成（在程序性协议中称为冻结后评估器）是一个多核CNN，具有并行的3、7、15和21 nt分支、逐点融合、深度卷积以及独立的均值和方差头~\citep{He2016,Lakshminarayanan2017}。使用种子集成和预测方差遵循已建立的神经网络不确定性策略~\citep{Gal2016,Blundell2015,Lakshminarayanan2017}。
每种架构使用三个固定种子：评审器为20260721--20260723，跨模型评估集成为20260731--20260733。

模型使用AdamW~\citep{Loshchilov2019}训练最多20个epoch，采用余弦学习率衰减、批次大小256、学习率$10^{-3}$、权重衰减$10^{-4}$、dropout 0.25~\citep{Srivastava2014}以及梯度裁剪5.0。损失结合了异方差高斯负对数似然、报告的测量标准误以及权重为0.1的反向互补一致性惩罚。早停使用验证集平均Pearson相关系数，耐心为6。测试记录未用于选择检查点、架构、超参数或方差尺度。

对于每个集成，预测是三个种子特定均值的平均，遵循自助聚合集成预测的精神~\citep{Breiman2001}。预测方差结合了平均学习到的模型内方差和种子间方差。每种细胞类型的一个乘法方差尺度仅从验证集的平方误差估计，然后冻结~\citep{Ovadia2019}。跨模型评估集成从未被搜索代码导入；它仅在候选文件及其SHA-256哈希写入后（程序性冻结）应用，确保生成与评估完全分离。

\subsection{编辑方法}

所有方法均针对靶细胞K562、HepG2和SK-N-SH，在精确替换预算1、5、10和20下进行评估。沿搜索路径每个位置最多编辑一次。

\paragraph{随机匹配编辑。}
位置和非参考等位基因均匀无放回抽样。600亲本评估使用三次独立固定种子（20260801--20260803）。方法级比较时，数值终点在重复间取平均；从不选择最优随机重复。

\paragraph{贪婪Malinois。}
每一步，每个可用的单核苷酸替换都由主特异性裕度打分，保留得分最高的替换。搜索期间不执行评审、不确定性或序列域约束。

\paragraph{主束（计算量匹配对照）。}
一束24个部分编辑集每次扩展一个替换。每步最多保留96个候选，按主得分排序——主得分为主裕度减去0.05倍正向/反向链不一致性。束宽、最大扩展深度和预筛选大小与SafeEdit-CRE完全匹配。不使用评审得分、不确定性惩罚或序列域重排序。

\paragraph{SafeEdit-CRE。}
SafeEdit-CRE使用与主束相同的主扩展和计算预算。主候选的佼佼者由评审集成根据下式重新打分
\begin{equation}
S(x)=M_t^{\mathrm{primary}}(x)+0.3M_t^{\mathrm{reviewer}}(x)
-0.1D_{\mathrm{strand}}(x)-0.05U_{\mathrm{reviewer}}(x).
\end{equation}
候选经过预过滤，检查同源多聚体长度是否超过8、以及相对于亲本的绝对GC变化是否超过0.12。重排序使用验证导出的、预算特定的链不一致性、评审不确定性、6聚体天然性变化、GC变化和同源多聚体长度限值。优先选择正的主裕度和评审裕度；当没有候选通过每个首选条件时，回退得分保留一条路径。在每个请求的预算处独立取快照。

序列天然性是重叠6聚体在天然训练序列频率下的位置无关平均对数似然。它在重排序期间用作计算序列域诊断。

\subsection{九准则设计基准}

九准则基准（G4）使用90条测试集CRE亲本，与开发或主要评估亲本无重叠（96个名义亲本中有6个在发现与探索性基准的标识符级重叠后被排除，未参考结果）。对于每个方法--亲本--靶标--预算组合，记录九项布尔检查：非负主靶标增益、正主裕度增益、正评审裕度传递、链不一致性在域内、评审不确定性在域内、天然性在域内、GC变化在域内、同源多聚体安全性、精确编辑预算。阈值从验证集随机编辑校准，并在评估前锁定。不可行组合必须保留在结果表中并标记为失败。

主要对比是SafeEdit-CRE与贪婪之间的配对九准则通过率差值。我们使用100,000次重复的亲本簇自助法计算95\%置信区间，保留每个亲本贡献的12个靶标--预算条件之间的依赖关系。保留汇总二元结果上的精确双侧McNemar检验作为描述性敏感性分析。使用分层的靶标--预算汇总来检验异质性。

\subsection{主要跨模型评估（600亲本）}

主要评估（G8）评估了600个标签盲、先前未使用的亲本，方法包括随机匹配、贪婪、主束和SafeEdit-CRE。搜索产生了43,200条序列行，因为随机编辑有三次重复。方法级统计时，三次随机重复平均为每个亲本--靶标--预算组合一个数值记录，共28,800行。此折叠表仅用于数值比较，从不用于选择DNA序列。候选文件在跨模型集成推理前冻结（写入SHA-256哈希）。

主要终点是跨模型特异性裕度增益。次要终点是跨模型靶标增益、正裕度比例和跨模型集成不确定性。对于每个对比，指标首先在一个亲本内的12个靶标--预算条件上平均；然后100,000次自助采样对600个亲本有放回重采样。不可行的束搜索结果作为零编辑、零增益的亲本序列保留。

\subsection{完整重搜索消融}

在180个亲本上独立重运行七次搜索：主束、完整SafeEdit-CRE、无评审得分的SafeEdit-CRE、无不确定性惩罚、无链惩罚、无天然性重排序、以及无序列预过滤器。与事后删除过滤器不同，每个消融都可能生成不同的序列。然后跨模型评估集成在全部15,120条候选行上重新运行。我们验证了候选键、精确编辑预算、有限数值、以及跨模型靶标和裕度增益的独立计算。配对亲本簇自助区间使用与主要评估相同的程序。

\subsection{候选分层}

3,240条九准则基准候选被组织为三个计算优先级层级。Tier A要求：SafeEdit-CRE生成、通过全部九项准则、非负靶标增益、正主裕度和评审裕度增益、精确编辑预算、无合成风险标记、且位于或接近双目标主--评审帕累托前沿。近帕累托候选指在前沿的5\%相对或0.1绝对范围内。Tier B通过核心活性和跨模型检查，但允许帕累托次优或轻度天然性偏差。Tier C未通过一项或多项硬性检查，仅保留用于诊断分析。补充表~S10提供了跨越K562、HepG2和SK-N-SH的三个代表性Tier A示例。

\subsection{软件、可重复性与报告}

初步开发阶段在NVIDIA RTX 3060笔记本GPU上运行。最终基准在Slurm管理的RTX 3080 Ti计算节点上运行。固定种子、环境记录、代码哈希、候选哈希、检查点哈希和运行日志均保存在冻结发布包中。统计值和图件由独立分析脚本直接从候选级表格重新生成，使用PyTorch、NumPy、pandas和Matplotlib实现~\citep{Paszke2019,Harris2020,McKinney2010,Hunter2007}。全部代码、环境文件、最小复现脚本、核心统计输入和候选序列表均包含在v1.1.0发布包中。
